\chapter{الحساب: القواسم والأعداد الأولية}\label{ch:g9:arith}

يدرس الحساب الأعداد الصحيحة وكيف \omterm{def:g9:arith:divisor}{يقسم} بعضها بعضًا. وأبطاله
الرئيسيون هم \omterm{def:g9:arith:prime}{الأعداد الأولية}، اللبنات التي يُبنى منها كل
عدد صحيح بالضرب. وينتهي الفصل عند \omterm{def:g9:arith:gcd}{القاسم المشترك الأكبر}، وهو الأداة
المناسبة لِاختزال الكسور مرّة واحدة وإلى الأبد. وتتواصل هذه
الحكاية، إلى أبعد من ذلك بكثير، في مجلّد الثانوية وما بعده.

\section{القواسم والمضاعفات}

\begin{definition}[القاسم، المضاعف]\label{def:g9:arith:divisor}
ليكن $a$ و $b$ عددين صحيحين موجبين. نقول إنّ $b$
\emph{يقسم}\index{قاسم} $a$ (أو إنّ $b$ قاسم للعدد $a$، أو إنّ
$a$ \emph{مضاعف}\index{مضاعف} للعدد $b$) حين يكون $a = b \times
k$ من أجل عدد صحيح $k$ --- أي حين تترك قسمة $a$ على $b$
\omterm{def:g3:division:remainder}{الباقي} $0$.
\end{definition}

\begin{example}
قواسم العدد $24$ هي $1, 2, 3, 4, 6, 8, 12, 24$ --- وهي تأتي
أزواجًا \omterm{def:g2:mult:def}{جداء} كلٍّ منها $24$: $(1,24)$، $(2,12)$، $(3,8)$، $(4,6)$.
\omterm{def:g5:division:multiple}{ومضاعفات} العدد $7$ هي $7, 14, 21, 28, \dots$
\end{example}

\begin{proposition}[قواعد قابلية القسمة]\label{prop:g9:arith:rules}
العدد الصحيح \omterm{def:g6:wholes:divisible}{قابل للقسمة}:
\begin{itemize}
  \item على $2$ حين يكون رقمه الأخير زوجيًا ($0, 2, 4, 6, 8$)؛
  \item على $5$ حين يكون رقمه الأخير $0$ أو $5$؛
  \item على $10$ حين يكون رقمه الأخير $0$؛
  \item على $3$ (وكذلك $9$) حين يكون مجموع أرقامه \omterm{def:g6:wholes:divisible}{قابلًا للقسمة}
        على $3$ (وكذلك $9$)؛
  \item على $4$ حين يشكّل رقماه الأخيران عددًا \omterm{def:g6:wholes:divisible}{قابلًا للقسمة} على $4$.
\end{itemize}
\end{proposition}
\admitted

\begin{example}
العدد $7\,215$ ينتهي بالرقم $5$: فهو \omterm{def:g6:wholes:divisible}{قابل للقسمة} على $5$. \omterm{def:g1:addition:def}{ومجموع}
أرقامه $7 + 2 + 1 + 5 = 15$، وهو \omterm{def:g6:wholes:divisible}{قابل للقسمة} على $3$ لا على $9$:
إذن $7\,215$ \omterm{def:g6:wholes:divisible}{قابل للقسمة} على $3$ لا على $9$. وبالفعل
$7\,215 = 3 \times 5 \times 481$.
\end{example}

\section{الأعداد الأولية}

\begin{definition}[العدد الأولي]\label{def:g9:arith:prime}
\emph{العدد الأولي}\index{عدد أولي} عدد صحيح $\geq 2$
قواسمه الوحيدة هي $1$ ونفسه. والأعداد الأولية الأصغر من $30$ هي
\[
2,\ 3,\ 5,\ 7,\ 11,\ 13,\ 17,\ 19,\ 23,\ 29 .
\]
والعدد $1$ \emph{ليس} أوليًا (باصطلاح)، وكل عدد صحيح $\geq 2$
غير أولي يُسمّى \emph{مركّبًا}.
\end{definition}

\begin{theorem}[التفكيك إلى عوامل أولية]\label{thm:g9:arith:factorization}
كل عدد صحيح $\geq 2$ هو \omterm{def:g2:mult:def}{جداء} \omterm{def:g9:arith:prime}{أعداد أولية}، وهذا
التفكيك وحيد إلى حدّ ترتيب العوامل.
\end{theorem}
\admitted

\begin{method}[تفكيك عدد صحيح]\label{met:g9:arith:factorization}
اِقسم على أصغر \omterm{def:g9:arith:prime}{عدد أولي} ممكن، مرارًا وتكرارًا، حتى تصل إلى $1$:
\begin{enumerate}
  \item جرّب $2$ ما دام العدد زوجيًا؛
  \item ثم جرّب $3$، ثم $5$، ثم $7$، \dots\ (\omterm{def:g9:arith:prime}{الأعداد الأولية} وحدها)؛
  \item وقف حين يصير \omterm{def:g3:division:remainder}{خارج القسمة} $1$؛ واِجمع العوامل مع
        أُسُسها.
\end{enumerate}
ويكفي أن تجرّب \omterm{def:g9:arith:prime}{الأعداد الأولية} $p$ التي لا يتجاوز $p^2$ العدد
الحالي: فإن لم يقسمه أيٌّ منها، كان العدد نفسه أوليًا.
\end{method}

\begin{example}\label{ex:g9:arith:factorization}
فكّك العدد $360$، قسمةً قسمة:
\[
360 = 2 \times 180, \quad
180 = 2 \times 90, \quad
90 = 2 \times 45, \quad
45 = 3 \times 15, \quad
15 = 3 \times 5,
\]
إذن
\[
360 = 2 \times 2 \times 2 \times 3 \times 3 \times 5 = 2^3 \times 3^2
\times 5 .
\]
\end{example}

\begin{omfigure}
\begin{tikzpicture}[scale=0.78]
  \node (n360) at (0,0) {$360$};
  \node (n180) at (1.2,-1) {$180$};
  \node (n90) at (2.4,-2) {$90$};
  \node (n45) at (3.6,-3) {$45$};
  \node (n15) at (4.8,-4) {$15$};
  \node (n5) at (6.0,-5) {$5$};
  \node[omThm] (p2a) at (-1.2,-1) {$2$};
  \node[omThm] (p2b) at (0,-2) {$2$};
  \node[omThm] (p2c) at (1.2,-3) {$2$};
  \node[omThm] (p3a) at (2.4,-4) {$3$};
  \node[omThm] (p3b) at (3.6,-5) {$3$};
  \draw (n360) -- (n180); \draw (n360) -- (p2a);
  \draw (n180) -- (n90); \draw (n180) -- (p2b);
  \draw (n90) -- (n45); \draw (n90) -- (p2c);
  \draw (n45) -- (n15); \draw (n45) -- (p3a);
  \draw (n15) -- (n5); \draw (n15) -- (p3b);
\end{tikzpicture}

{\small شجرة عوامل العدد $360$: كل خطوة تفصل أصغر عامل أولي
(بالأحمر). وبقراءة الأوراق الحمراء والعدد الأخير $5$:
$360 = 2^3 \times 3^2 \times 5$.}
\end{omfigure}

\begin{theorem}[إقليدس]\label{thm:g9:arith:euclid}
\omterm{def:g9:arith:prime}{الأعداد الأولية} غير منتهية العدد.
\end{theorem}

\begin{proof}
لنفترض أنّها منتهية العدد، ولتكن $p_1, p_2, \dots, p_k$، ولنتأمّل
العدد
\[
N = p_1 \times p_2 \times \dots \times p_k + 1 .
\]
قسمة $N$ على أيّ $p_i$ تترك \omterm{def:g3:division:remainder}{الباقي} $1$، فلا \omterm{def:g9:arith:divisor}{يقسم} أيٌّ من
الأعداد $p_i$ العدد $N$. لكن $N \geq 2$ له قاسم أولي واحد على
الأقل (\cref{thm:g9:arith:factorization}) --- وهو \omterm{def:g9:arith:prime}{عدد أولي} ليس في
قائمتنا. وهذا تناقض: فلا قائمة منتهية تسع كل \omterm{def:g9:arith:prime}{الأعداد الأولية}.
\end{proof}

\section{القاسم المشترك الأكبر}

\begin{definition}[القاسم المشترك الأكبر]\label{def:g9:arith:gcd}
\emph{القاسم المشترك الأكبر}\index{قاسم مشترك أكبر} لعددين صحيحين
موجبين $a$ و $b$، ويُكتب $\gcd(a, b)$، هو أكبر عدد صحيح يقسمهما
معًا. وحين يكون $\gcd(a, b) = 1$، يُقال إنّ العددين \emph{أوليان
فيما بينهما}\index{أوليان فيما بينهما}: فلا يشتركان في أيّ
قاسم سوى $1$.
\end{definition}

\begin{example}
قواسم العدد $18$: $1, 2, 3, 6, 9, 18$. وقواسم العدد $24$:
$1, 2, 3, 4, 6, 8, 12, 24$. والقواسم المشتركة: $1, 2, 3, 6$؛ إذن
$\gcd(18, 24) = 6$. والعددان $15$ و $28$ \omterm{def:g9:arith:gcd}{أوليان فيما بينهما}.
\end{example}

\begin{proposition}[القاسم المشترك الأكبر من التفكيكين]\label{prop:g9:arith:gcdfact}
\omterm{def:g9:arith:gcd}{القاسم المشترك الأكبر} لعددين صحيحين هو \omterm{def:g2:mult:def}{جداء} \omterm{def:g9:arith:prime}{الأعداد الأولية}
التي تظهر في \emph{كلا} التفكيكين، مع أخذ \emph{الأصغر} من
\omterm{def:g8:powers:def}{أُسّيه} في كل مرّة.
\end{proposition}
\admitted

\begin{example}\label{ex:g9:arith:gcdfact}
$360 = 2^3 \times 3^2 \times 5$ و $84 = 2^2 \times 3 \times 7$.
\omterm{def:g9:arith:prime}{والأعداد الأولية} المشتركة: $2$ (\omterm{def:g8:powers:def}{الأُسّان} $3$ و $2$: نُبقي $2$)
و $3$ (\omterm{def:g8:powers:def}{الأُسّان} $2$ و $1$: نُبقي $1$). إذن
\[
\gcd(360, 84) = 2^2 \times 3 = 12 .
\]
\end{example}

\begin{theorem}[خوارزمية إقليدس]\label{thm:g9:arith:euclidalgo}
إذا كانت $a = bq + r$ قسمة $a$ على $b$ حيث \omterm{def:g3:division:remainder}{الباقي} هو $r$، فإنّ
\[
\gcd(a, b) = \gcd(b, r).
\]
وبتكرار القسمات حتى يصير \omterm{def:g3:division:remainder}{الباقي} $0$، يكون \omterm{def:g9:arith:gcd}{القاسم المشترك الأكبر}
للعددين $a$ و $b$ هو \emph{آخر باقٍ غير معدوم}.
\end{theorem}

\begin{proof}
من $a = bq + r$: كل عدد صحيح \omterm{def:g9:arith:divisor}{يقسم} $b$ و $r$ \omterm{def:g9:arith:divisor}{يقسم} $bq + r = a$؛
ومن $r = a - bq$: كل عدد صحيح \omterm{def:g9:arith:divisor}{يقسم} $a$ و $b$ \omterm{def:g9:arith:divisor}{يقسم} $r$. إذن للزوجين
$(a, b)$ و $(b, r)$ القواسم المشتركة نفسها تمامًا --- وبخاصّة
أكبرها نفسه. وبما أنّ البواقي تتناقص تناقصًا تامًّا، تنتهي
الخوارزمية، ويعطي $\gcd(x, 0) = x$ آخر باقٍ غير معدوم.
\end{proof}

\begin{example}\label{ex:g9:arith:euclidalgo}
اِحسب $\gcd(1071, 462)$:
\begin{align*}
1071 &= 462 \times 2 + 147, \\
462 &= 147 \times 3 + 21, \\
147 &= 21 \times 7 + 0 .
\end{align*}
آخر باقٍ غير معدوم هو $21$: أي $\gcd(1071, 462) = 21$.
\end{example}

\begin{method}[اِختزال كسر اِختزالًا تامًّا]\label{met:g9:arith:simplify}
لكتابة $\dfrac ab$ على أبسط صورة:
\begin{enumerate}
  \item اِحسب $d = \gcd(a, b)$، بخوارزمية إقليدس مثلًا؛
  \item واِقسم \omterm{def:g4:fractions:def}{البسط} \omterm{def:g4:fractions:def}{والمقام} على $d$:
        $\dfrac ab = \dfrac{a \div d}{b \div d}$؛
  \item فيكون \omterm{def:g9:fractions:fraction}{الكسر} الناتج \emph{غير قابل للاِختزال}: بسطه
        ومقامه \omterm{def:g9:arith:gcd}{أوليان فيما بينهما}.
\end{enumerate}
\end{method}

\begin{example}
$\dfrac{462}{1071} = \dfrac{462 \div 21}{1071 \div 21} = \dfrac{22}{51}$،
و $\gcd(22, 51) = 1$: فهو غير قابل للاِختزال.
\end{example}

\section{تمارين}

\begin{exercise}[$\star$]\label{exo:g9:arith:1}
اُذكر جميع قواسم العدد $36$، والعدد $45$، والعدد $17$.
\end{exercise}

\begin{exercise}[$\star$]\label{exo:g9:arith:2}
باِستعمال قواعد \omterm{def:g6:wholes:divisible}{قابلية القسمة}، حدّد ما إذا كان $2\,346$
\omterm{def:g6:wholes:divisible}{قابلًا للقسمة} على $2$، و $3$، و $4$، و $5$، و $9$.
\end{exercise}

\begin{exercise}[$\star$]\label{exo:g9:arith:3}
أعطِ التفكيك إلى عوامل أولية للأعداد $72$ و $150$ و $210$ و $121$.
\end{exercise}

\begin{exercise}[$\star$]\label{exo:g9:arith:4}
هل $101$ أولي؟ وهل $91$؟ وهل $143$؟ برّر باِستعمال قاعدة التوقف
في \cref{met:g9:arith:factorization}.
\end{exercise}

\begin{exercise}[$\star$]\label{exo:g9:arith:5}
اِحسب $\gcd(48, 60)$ بطريقتين: بسرد القواسم المشتركة، ثم
اِنطلاقًا من التفكيكين إلى عوامل أولية.
\end{exercise}

\begin{exercise}[$\star\star$]\label{exo:g9:arith:6}
اِستعمل خوارزمية إقليدس لحساب $\gcd(255, 154)$، ثم
$\gcd(1053, 325)$. واُكتب كل سطر من سطور القسمة.
\end{exercise}

\begin{exercise}[$\star\star$]\label{exo:g9:arith:7}
اِجعل \omterm{def:g9:fractions:fraction}{الكسر} $\dfrac{588}{504}$ غير قابل للاِختزال. (اِحسب \omterm{def:g9:arith:gcd}{القاسم
المشترك الأكبر} بالطريقة التي تختارها، ثم اِقسم.)
\end{exercise}

\begin{exercise}[$\star\star$]\label{exo:g9:arith:8}
لدى بائعة زهور $84$ وردة و $126$ زهرة توليب، وتريد أن تصنع باقات
متطابقة تستعمل فيها كل الزهور، وأن يكون عدد الباقات أكبر ما يمكن.
كم باقة تستطيع أن تصنع؟ وماذا تحوي كل واحدة؟
\end{exercise}

\begin{exercise}[$\star\star$]\label{exo:g9:arith:9}
تنطلق عبّارتان من الرصيف نفسه في الساعة 8:00. تنطلق إحداهما كل
$24$ دقيقة، والأخرى كل $36$ دقيقة. في أيّ ساعة تنطلقان معًا في
المرّة التالية؟ (اِبحث عن أصغر \omterm{def:g9:arith:divisor}{مضاعف} مشترك للعددين $24$ و $36$؛
والتفكيكان يساعدان.)
\end{exercise}

\begin{exercise}[$\star\star\star$]\label{exo:g9:arith:10}
ليكن $n$ عددًا صحيحًا موجبًا.
\begin{enumerate}
  \item بيّن أنّ $\gcd(n, n+1) = 1$ (العددان الصحيحان المتتاليان
        \omterm{def:g9:arith:gcd}{أوليان فيما بينهما} دائمًا).
  \item واِستنتج أنّ \omterm{def:g9:fractions:fraction}{الكسر} $\dfrac{n}{n+1}$ غير قابل للاِختزال دائمًا.
\end{enumerate}
\end{exercise}

\section{مسألة: أباريق الماء، والزيزان، ومئة خزانة}

\begin{problem}[{مسألة نهاية الأسبوع --- القاسم المشترك الأكبر
يقرّر أيّ المقادير يمكن قياسها بإبريقين؛ والأعداد الأولية تحمي
الزيزان؛ والخزانات التي تبقى مفتوحة هي المربّعات التامّة}]
\label{pb:g9:arith:1}
ثلاثة ألغاز تبدو أحاجي وهي في الحقيقة حساب: قياس الماء بأوانٍ
بلا تدريج (\omterm{def:g9:arith:gcd}{القاسم المشترك الأكبر} متنكّرًا)، ودورات حياة حشرات
تطوّرت إلى \omterm{def:g9:arith:prime}{أعداد أولية}، ورواق شهير فيه مئة خزانة تتقرّر
حالته النهائية بعدّ القواسم. وكل شيء هنا يشتغل بعُدّة هذا
الفصل: \omterm{def:g6:wholes:divisible}{قابلية القسمة}، والتفكيك إلى عوامل أولية
(\cref{thm:g9:arith:factorization})، وخوارزمية إقليدس
(\cref{thm:g9:arith:euclidalgo}).

\medskip
\noindent\textbf{الجزء الأول --- أباريق الماء.}
أنت واقف عند نافورة ومعك إبريقان بلا تدريج، سعتهما $5$ لترات
و $3$ لترات. والحركات المسموحة: ملء إبريق إلى حافّته، وتفريغ
إبريق تفريغًا تامًّا، وصبّ إبريق في الآخر حتى يفرغ المصدر أو
يمتلئ الهدف.
\begin{enumerate}
  \item قِس $1$ لتر بالضبط. (صِف تسلسل حركاتك ومحتوى الإبريقين
        بعد كل حركة.)
  \item وقِس $4$ لترات بالضبط --- وهو اللغز الوارد في فيلم حركة
        شهير. (يمكن ذلك في ست حركات.)
  \item وأيّ أعداد صحيحة من اللترات بين $1$ و $8$ تستطيع أن
        تُظهر (في إبريق واحد، أو موزّعة على الاِثنين)؟ أكمل
        القائمة معيدًا اِستعمال تسلسلاتك.
  \item وإبريقان جديدان: $6$ لترات و $4$ لترات. حاول قياس $1$
        لتر --- ثم فسّر لماذا الأمر ميؤوس منه: تحقّق أنّ كل حركة
        من الحركات الثلاث المسموحة تُبقي محتوى كل إبريق
        مضاعفًا للعدد $2$، فيكون كل مقدار يمكن بلوغه زوجيًا.
  \item وحجّة السؤال 4 تصلح في العموم: بإبريقين سعتهما $a$ و $b$
        لترات، يكون كل مقدار يمكن بلوغه مضاعفًا للعدد
        $\gcd(a, b)$. اِحسب $\gcd(6, 4)$ و $\gcd(5, 3)$، وقل ماذا
        يتنبّأ به القانون في كل زوج من الأباريق.
\end{enumerate}

\medskip
\noindent\textbf{الجزء الثاني --- إقليدس عند النافورة.}
\begin{enumerate}[resume]
  \item اِحسب بخوارزمية إقليدس: $\gcd(91, 65)$ و $\gcd(2\,026, 46)$.
  \item وفسّر بكلماتك لماذا تكون المقادير الظاهرة في الأباريق
        بواقي إقليدس متنكّرة: بإبريقين سعتهما $13$ لترًا
        و $5$ لترات، اِملأ الإبريق الصغير مرارًا واِصببه في
        الكبير (مفرّغًا الكبير كلّما اِمتلأ). ما المقادير الجديدة
        التي تظهر أولًا --- وقارنها مع البواقي في خوارزمية
        إقليدس من أجل $(13, 5)$.
  \item واِستنتج جواب البطل: بإبريقين سعتهما $13$ و $5$ لترات، هل
        تستطيع قياس $1$ لتر بالضبط؟ برّر في سطر واحد بالسؤال 5
        و $\gcd(13, 5)$.
  \item وبرهان سريع على الأولية فيما بينهما على غرار
        \cref{exo:g9:arith:10}: بيّن أنّ $\gcd(n, 2n + 1) = 1$ من
        أجل كل عدد صحيح موجب $n$. (ما الذي يجب أن يقسمه
        قاسم مشترك للعددين $n$ و $2n + 1$؟)
  \item وتنطلق حافلتان معًا من المحطّة النهائية في الساعة 7:00؛
        تنطلق إحداهما كل $12$ دقيقة، والأخرى كل $18$. اُذكر أوقات
        اِنطلاق كل منهما التالية، وجِد أول لحظة تنطلقان فيها معًا
        من جديد. وتحقّق في هذا المثال من القانون الجميل: (أول
        \omterm{def:g9:arith:divisor}{مضاعف} مشترك) $\times$ $\gcd$ $=$ \omterm{def:g2:mult:def}{جداء} العددين ---
        ثم اِختبره من جديد على $5$ و $3$.
\end{enumerate}

\medskip
\noindent\textbf{الجزء الثالث --- الزيزان والقواسم والخزانات.}
\begin{enumerate}[resume]
  \item تخرج بعض زيزان أمريكا الشمالية كل $17$ سنة فقط؛ ولنفترض
        أنّ عدد حيوان مفترس يبلغ ذروته كل $4$ سنوات. فإذا وقع
        الأمران هذه السنة، فبعد كم سنة يتصادف الخروج مع الذروة
        من جديد؟ والسؤال نفسه لو كانت دورة الزيزان $16$ سنة ---
        كل كم تُذبح حينئذ؟ وفسّر في جملة واحدة لماذا دفع التطوّر
        الدورة إلى طول \emph{أولي}.
  \item وباِستعمال التفكيك $360 = 2^3 \times 3^2 \times 5$، عُدّ
        قواسم العدد $360$ من غير أن تسردها: القاسم يختار
        \omterm{def:g8:powers:def}{أُسًّا} للعدد $2$ (أربعة اِختيارات: $0, 1, 2, 3$)،
        \omterm{def:g8:powers:def}{وأُسًّا} للعدد $3$، \omterm{def:g8:powers:def}{وأُسًّا} للعدد $5$. فكم قاسمًا في
        \omterm{def:g1:addition:def}{المجموع}؟
  \item وبيّن أنّه في تفكيك مربّع تامّ $n = m^2$، يكون كل \omterm{def:g9:arith:prime}{عدد
        أولي} \emph{زوجيّ} \omterm{def:g8:powers:def}{الأُسّ}. واِستنتج، من غير حساب
        أيّ جذر تربيعي، أنّ $360$ ليس مربّعًا تامًّا.
  \item وقارِن كل قاسم $d$ للعدد $n$ بشريكه $\frac{n}{d}$ (من
        أجل $n = 36$: $1 \leftrightarrow 36$، $2 \leftrightarrow 18$،
        $3 \leftrightarrow 12$، $4 \leftrightarrow 9$،
        $6 \leftrightarrow 6$). متى يكون القاسم شريك نفسه؟
        واِستنتج المعيار: يكون للعدد $n$ عدد \emph{فردي} من
        القواسم بالضبط حين يكون $n$ مربّعًا تامًّا. وتحقّق من
        ذلك على $36$ وعلى $360$.
  \item الخزانات المئة. الخزانات من $1$ إلى $100$ مغلقة في
        البداية. التلميذ $1$ يقلب حالة كل خزانة؛ والتلميذ $2$
        يقلب الخزانات $2, 4, 6, \dots$؛ والتلميذ $k$ يقلب
        \omterm{def:g5:division:multiple}{مضاعفات} $k$؛ وهكذا حتى التلميذ $100$. فسّر أيّ
        التلاميذ يمسّون الخزانة $n$، وكم مرّة تُقلب حالتها،
        و --- باِستعمال السؤال 14 --- أيّ الخزانات بالضبط تنتهي
        مفتوحة. وكم واحدة منها مفتوحة؟
\end{enumerate}
\end{problem}
